\documentclass{jsarticle}
\usepackage{amssymb}
\usepackage{amsthm}
\usepackage{amsmath}
\usepackage{comment}
%定理環境 thmはあまり使わずpropをなるべく使う。
%主張は基本的に証明の中で使い、そのため番号を付けない。
%注意環境を付けてみた。
\newtheorem{thm}{定理}
\newtheorem{prop}[thm]{命題}
\newtheorem{cor}[thm]{系}
\newtheorem{lem}[thm]{補題}
\newtheorem{df}[thm]{定義}
\newtheorem{exam}[thm]{例}
\newtheorem{fact}[thm]{事実}
\newtheorem*{claim}{主張}
\newtheorem*{cau}{注意}
\renewcommand{\proofname}{証明}
%矢印たち。
%\toは\rightarrowと同じ。\getsは\leftarrowと同じ。同値は\iff
\newcommand{\To}{\Longrightarrow}
\newcommand{\Gets}{\LongLeftarrow}
%黒板太字
\newcommand{\nn}{\mathbb{N}}
\newcommand{\zz}{\mathbb{Z}}
\newcommand{\qq}{\mathbb{Q}}
\newcommand{\rr}{\mathbb{R}}
\newcommand{\cc}{\mathbb{C}}
%無限大
\newcommand{\mgn}{\infty}
%差集合は\setminusで統一する。等号の否定は\neq
%真部分集合は\subsetneqq　偏微分は\partial
%ディスプレイスタイル。\dfracでディスプレイ分数
\newcommand{\dpst}{\displaystyle}

\title{$T_1$を仮定しない正規性が正則性を導出しない事を示す例}
\author{ナンブキトラ}
\date{\today}
\begin{document}
\maketitle
この文章では位相空間の正規性が正則性を導かないことを示す。$T_1$を課した意味での正規生からはもちろん正則性を導けるのだが、そうでない時は正規性と正則性は独立した概念なのである。筆者が学部２年生の時の時にそのような考えようとしてその時は見つけられなかった覚えがあるが、ある時ふと思い付いたのでここに記しておく。\\ 
以下の文中では正規、正則という言葉は$T_1$を積極的には仮定しない意味で用いて$T_3,T_4$は$T_1$を仮定した意味で使う。著者がこの言葉使いを採用した理由は定理2の見た目を良くするためであり、他の文章を読む時は定義をしかりと確認する必要がある。

\begin{df}分離公理位相空間$X$について以下の条件を定義する。\\
$(T_0)$\  異なる２点$x,y$について$x$の近傍で$y$を含まないようなものが存在するか、もしくは$y$の近傍で$x$を含まないようなものが存在する\\ 
$(T_1)$\ 異なる２点$x,y$について$x$の近傍で$y$を含まないようなものが存在する。\\
$(T_2$,もしくはハウスドルフ$)$\ 異なる２点$x,y$について$ｘ$の近傍$U,y$の近傍Vで$U\cap V=\O $ となるものが存在する。\\
$($正則$)$\ 互いに素なる一点$x$と閉集合$F$について$x$を含む開集合$U,F\subset V$となる開集合$V$で$U\cap V=\O$ となるものが存在する。\\
$($正規$)$\ 互いに素な２つの閉集合が開集合で分離される\\
$(T_3)$正則かつ$T_{1}$\\
$(T_{4})$正規かつ$T_{1}$\\
\end{df}
簡単にわかることなので注意にとどめるが、正則性と、各点の閉近傍全体がその点の基本近傍系を成すことは同値である。この事はこれから自在に用いる。また任意の空でない閉集合が交わるような空間は正規となることに注意しよう（所謂、空なる前提というやつである。）

\begin{fact}\label{fact}
\[
T_{4}\Rightarrow T_{3}\Rightarrow T_{2}\Rightarrow T_{1}\Rightarrow T_0
\]
\end{fact}

\begin{thm}[本題に関係ないおまけ]\label{omake}
正則かつ$(T_0)$ならば$(T_1)$も充たす。すなわちそのような空間はハウスドルフである。
\end{thm}
\begin{proof}ハウスドルフ性を示す。異なる２点$x,y$に対してそのいづれかの近傍が存在して他方の天を含まないが、$x$の近傍$U$が存在して$y\not\in U$と仮定しても一般性を失わない。さて正則性から$x$の閉近傍$F$が存在して$x\in F\subset U$となる。よって特に$y\not\in F$である$F$が閉であることと再び正則性から$F\subset V, y\in W,V\cap W=\emptyset$となる開集合$V,W$が存在する。この$V,W$が$x,y$を分離する開集合になるのでこの空間はハウスドルフである。
\end{proof}
さて本題に入ろう。事実\ref{fact}に対して次の疑問が起こる
\[
\text{（正規$\Rightarrow$正則）?}
\]
これは一般に成り立たない事を例を以って示そう。

\begin{exam}シェルピンスキー空間\\
$A=\{0,1\}$に
\[
\{\emptyset , \{0\},A\}
\]
を位相として入れた空間である。これが位相を定めることを確かめるのは容易い。\\ 
この空間の閉集合は
\[
\{\emptyset , \{1\},A\}
\]
の３つである。この事から$T_1$ではないことと、全ての空でない閉集合が交わるので正規となる。
\\ また閉集合$\{1\}$を含む開集合は$A$しかないことから１点$0$と閉集合$\{1\}$を分離する開集合はないので正則ではない。ちなみにこの空間は$T_0$を充たす。
\end{exam}
このシェルピンスキー空間は任意の空でない閉集合が交わるような空間である。そのことから正規性を導いているが、より多くを求めて、任意の空でない閉集合が交わってはいない空間で正規かつ正則でない例をあげよう。上の例から引き続き$A=\{1,2\}$とする。
\begin{exam}正規かつ正則ではない例:１点とシェルピンスキー空間の位相的直和\\ 
位相的直和がわからない人のために明示すると台集合は
\[
B=\{ 1,2,3\}=A\cup \{3\}
\]
開集合系は
\[
\{ \emptyset , \{3\},\{1\},A,\{3\}\cup \{1\},\{3\}\cup A \}
\]	である。$A$がシエルピンスキー空間。$\{3\}$が１点という位相空間を務めている。位相的直和とは二つの空間を交わらないように"並べる"操作である\\ 
この空間が$T_1$を充たさないこと、正則でないことは上と同様である。\\ 
正規であることは交わらない閉集合の対が$\{ 3 \}$と$A$、もしくは$\{3\}$と$\{1\}$しかなく、これらは$\{3\}$と$A$と言う開集合で分離できるのでわかる。やはりこの空間も$T_0$を充たす。
\end{exam}
また上の二つの例は定理\ref{omake}の正則性を正規性に置き換えると結論がもはや成り立たないことも示している。

\end{document}